% TOP_PROBLEM: {"problemId": "problem.the-mckay-navarro-conjecture-for-finite-groups", "canonicalTitle": "The McKay-Navarro Conjecture for Finite Groups", "releaseRank": 495, "primaryDomain": "algebra_representation_category", "primaryDomainLabel": "Algebra, representation & category theory", "displayStatus": "open_with_solved_subcases", "releaseStatus": "open_with_solved_subcases"}
% !TeX program = lualatex
% ATTEMPT_STATUS: unresolved
% Model: GPT-6 Astra Ultra; continuation dated 27 September 2026.
\documentclass[11pt]{article}
\usepackage[margin=25mm]{geometry}
\usepackage{fontspec}
\setmainfont{Latin Modern Roman}
\usepackage{amsmath,amssymb,mathtools}
\usepackage{unicode-math}
\setmathfont{Latin Modern Math}
\usepackage{xurl,hyperref}
\hypersetup{colorlinks=true,urlcolor=blue}
\setcounter{secnumdepth}{0}
\setlength{\parindent}{0pt}
\setlength{\parskip}{5pt}
\setlength{\emergencystretch}{3em}
\begin{document}
\section*{\texorpdfstring{The McKay-Navarro Conjecture for Finite Groups}{The McKay-Navarro Conjecture for Finite Groups}}\label{495}
\subsection{Definitions and mathematical statement}
Let $G$ be finite, $p$ prime, and $P$ a Sylow $p$-subgroup. Write $\operatorname{Irr}_{p'}(G)$ for its irreducible complex characters of degree not divisible by $p$. In $\mathcal G=\operatorname{Gal}({\mathbb{Q}}(\zeta_{|G|})/{\mathbb{Q}})$, let $\mathcal H_p$ consist of automorphisms whose action on roots of unity of order prime to $p$ is raising to a common $p$-power. It acts by $\chi^\sigma(g)=\sigma(\chi(g))$.

The standard conjecture stated in the cited paper asks for an $\mathcal H_p$-equivariant bijection
\[
 \operatorname{Irr}_{p'}(G)\longleftrightarrow
 \operatorname{Irr}_{p'}(N_G(P)).
\]
Equivariance means that applying any $\sigma\in\mathcal H_p$ before or after the bijection has the same effect. It does not require equivariance under arbitrary Galois automorphisms. The source proves the prime-two case and states the all-prime conjecture. The catalogue's extra word blockwise refers to a distinct refinement; no precise blockwise target is supplied there, and that refinement is not silently identified with this global bijection.
\par\smallskip\textit{Formulation and concept references:} \href{https://arxiv.org/abs/2211.14237}{[S1]}.

\subsection{Short English statement}
Can prime-to-p-degree characters of a finite group and its Sylow normalizer be matched compatibly with the specified Galois action?
\subsection{Research attempt}

\paragraph{Exact target and source audit (27 September 2026).}
The target is an isomorphism of two finite $\mathcal H_p$-sets. Neither
equality of their cardinalities nor an arbitrary ordinary McKay bijection
is enough. The normalizer's characters are viewed in the same cyclotomic
field as the characters of $G$: its order divides $|G|$, so its character
values belong to that field. We retain the supplied global statement and
do not add a blockwise requirement.

The accepted author version of [S1], arXiv:2211.14237v2, explicitly states
the $\mathcal H_p$-equivariant conjecture and proves the prime-two case in
Corollary B [E1]. The newer manuscript \emph{The Isaacs--Navarro Galois
Conjecture}, arXiv:2509.02300v1, proves an equivariant bijection for the
subgroup $\mathcal H_0$ acting trivially on prime-to-$p$ roots and through
$p$-power-order automorphisms in finite cyclotomic quotients [E2,
Theorem A]. This is a substantial subcase, but it does not assert
equivariance for the full $\mathcal H_p$ in this entry. Its introduction
also records the completion of ordinary McKay. Those distinctions prevent
an incorrect inference that the all-prime target was settled by merely
counting characters or by the newer smaller-subgroup theorem.

There is a bibliographic discrepancy worth recording without changing the
supplied source entry: [S1] lists volume 483, whereas the institutional
publication record for the same article number 110369 gives
\emph{Advances in Mathematics} 477 [E3]. The mathematical formulation
was checked against the actual accepted author text [E1].

The attack has two parts. First convert equivariance into exact fixed-point
counting conditions and stress-test which counts are sufficient. Then
construct the bijection explicitly for a family with a nonnormal Sylow
subgroup, where the relevant characters and Galois action can both be
computed from first principles. The aim is to isolate the missing
information when that construction no longer applies.

\paragraph{Make the acting group finite and explicit.}
Write $n=|G|=p^a m$ with $(p,m)=1$. Each automorphism of
$\mathbb Q(\zeta_n)$ is $\sigma_u:\zeta_n\mapsto\zeta_n^u$ for a
unit $u\bmod n$. The condition defining $\mathcal H_p$ is exactly
\[
 u\bmod m\in\langle p\bmod m\rangle.
\]
To verify this equivalence, it suffices to check the primitive $m$th root:
if it is raised to $p^j$, the same exponent works on every root whose
order divides $m$. Conversely such a common power must work on that
primitive root. The action on $p$-power roots is unrestricted. The Chinese
remainder theorem therefore gives
\begin{equation}\label{ra495:H}
 \mathcal H_p\cong
 (\mathbb Z/p^a\mathbb Z)^\times\times
                  \langle p\bmod m\rangle.
\end{equation}
The factor of modulus one is interpreted as the trivial group. This
description uses roots with orders dividing the chosen cyclotomic
conductor; it does not require a prescribed action on roots outside that
field. The finite group $H=\mathcal H_p$ is abelian, but need not be
cyclic even when $p$ is odd.

For example, for $p=3,n=24$, the permitted exponents modulo $24$ are
\[
 1,11,17,19.
\]
Their reductions modulo $8$ lie in $\{1,3\}=\langle3\rangle$, and
their reductions modulo $3$ are arbitrary units. Every nonidentity
element has order two, so $H\cong C_2\times C_2$.
Thus an argument that silently treats $\mathcal H_p$ as generated by
one automorphism is not valid in the stated scope.

If $P$ is normal in $G$, then $N_G(P)=G$ and the identity character map
already proves the target. This includes $p\nmid|G|$, when $P=1$,
and all $p$-groups. The nontrivial part of the problem is the interaction
between character values and a proper Sylow normalizer.

\paragraph{An exact finite-set reduction using all subgroup marks.}
For a finite $H$-set $X$ and a subgroup $U\leq H$, write
\[
 F_X(U)=|X^U|=|\{x:\sigma x=x\text{ for every }\sigma\in U\}|.
\]
Call this a subgroup fixed-point count. Since $H$ is abelian, every
transitive orbit is isomorphic to $H/K$, and every point in that orbit
has stabilizer exactly $K$. Let $b_X(K)$ be the number of such orbits.
Then
\begin{equation}\label{ra495:marks}
 F_X(U)=\sum_{K\geq U} b_X(K)\,[H:K].
\end{equation}
Indeed $U$ fixes every coset of $K$ if $U\leq K$, and fixes none
otherwise; commutativity is used at this step.

\textbf{Lemma 495.1 (fixed-subgroup criterion).}
Two finite sets with an action of a finite abelian group $H$ admit an
$H$-equivariant bijection if and only if
$F_X(U)=F_Y(U)$ for every subgroup $U\leq H$.

\emph{Proof.}
A bijection preserves each fixed subset, so the condition is necessary.
For sufficiency, start at $K=H$ in \eqref{ra495:marks} and proceed by
decreasing subgroup order. The recursive formula
\[
 b_X(K)=\frac{F_X(K)-\sum_{J>K}b_X(J)[H:J]}{[H:K]}
\]
determines every orbit multiplicity uniquely. Equal fixed-point counts
give $b_X(K)=b_Y(K)$ for every $K$. Pair orbits with the same stabilizer,
choose one point in each paired orbit, and map $h x$ to $h y$.
This is well-defined because the stabilizers agree, and the resulting
disjoint union of orbit maps is an equivariant bijection.\hfill$\square$

For the present conjecture, the criterion says that it is enough, and
necessary, to prove
\begin{equation}\label{ra495:targetmarks}
 |\operatorname{Irr}_{p'}(G)^U|
   =|\operatorname{Irr}_{p'}(N_G(P))^U|
                 \qquad\text{for every }U\leq\mathcal H_p.
\end{equation}
The equality for $U=1$ is only the ordinary count. Equality for each
individual automorphism is the equality for cyclic subgroups. The
distinction between these two levels and all subgroups is mathematically
real, as the next construction shows.

\paragraph{A counterexample to an insufficient intermediate lemma.}
Let $H=C_2\times C_2$, with its three subgroups of order two denoted
$K_1,K_2,K_3$. Consider the $H$-sets
\[
 X=(H/1)\sqcup(H/H)\sqcup(H/H),\qquad
 Y=(H/K_1)\sqcup(H/K_2)\sqcup(H/K_3).
\]
Both have six elements. The identity fixes all six on either side.
A nonidentity element $h$ belongs to exactly one $K_i$, so on $Y$ it
fixes exactly the two points in $H/K_i$ and none in the other orbits.
On $X$ it fixes exactly the two singleton orbits. Thus
\[
 |X^h|=|Y^h|\qquad\text{for every }h\in H.
\]
Nevertheless $|X^H|=2$ and $|Y^H|=0$. There is no equivariant
bijection. Equivalently, equal permutation characters need not imply
isomorphic finite $H$-sets.

This does not construct a finite group violating McKay--Navarro: the two
sets have not been realized as the required character sets. It does
disprove the proposed general deduction from individual Galois fixed
counts to an equivariant bijection. Because the actual acting group can
be $C_2\times C_2$ as in \eqref{ra495:H}, that deduction cannot be used
without additional character-theoretic structure. In contrast, if $H$
is cyclic, every subgroup has a generator and individual fixed counts
do determine all the numbers in \eqref{ra495:targetmarks}. The failed
lemma is therefore precisely an issue for noncyclic actions, rather than
a universal failure of fixed-point counting.

\paragraph{A nontrivial family with a canonical explicit bijection.}
We now work with groups
\[
 G=A\rtimes P,
\]
where $A$ is a finite abelian group of order prime to $p$ and $P$ is a
finite abelian $p$-group acting on $A$. Faithfulness and fixed-point-free
action are not required. The Sylow subgroup $P$ need not be normal.
This is an elementary known-type subcase, proved here directly; no claim
of novelty or of an unrestricted reduction to this family is intended.

Write $A$ additively and put $C=C_A(P)$. Let $e=\operatorname{exp}(A)$,
and choose an integer $r$ such that $r|P|\equiv1\pmod e$. Define
\begin{equation}\label{ra495:average}
 T(a)=r\sum_{u\in P}u\cdot a.
\end{equation}
This is a homomorphism onto $C$. The sum is invariant under $P$ by
permutation of the terms, and for $c\in C$ it equals $r|P|c=c$.
Consequently $T^2=T$.

Let $B=[A,P]$ denote the subgroup generated by all $u\cdot a-a$.
The map $T$ kills each generator, so $B\leq\ker T$. In $A/B$ the
$P$ action is trivial, and therefore
$T(a)\equiv r|P|a\equiv a\pmod B$. If $T(a)=0$, this shows
$a\in B$. It follows that
\begin{equation}\label{ra495:decomposition}
 A=C\oplus B,\qquad B=\ker T.
\end{equation}
For completeness, each $a$ decomposes as $T(a)+(a-T(a))$, with the
second term in the kernel; the intersection is zero because $T$ is the
identity on $C$. This averaging works because the group order $|P|$ is
invertible modulo $e$. It is an actual endomorphism of the finite
abelian group, not an informal average of character values.

Let $\widehat A$ be the set of linear complex characters of $A$.
A character $\lambda\in\widehat A$ is $P$-invariant exactly when it
is trivial on $B$. Restriction therefore gives the bijection
\begin{equation}\label{ra495:dualrestriction}
 \widehat A^{P}\longrightarrow\widehat C,
 \qquad\lambda\longmapsto\lambda|_C,
\end{equation}
whose inverse sends $\eta$ to $\eta\circ T$. Invariance is equivalent
to triviality on each difference $u\cdot a-a$, and the decomposition
\eqref{ra495:decomposition} proves both injectivity and surjectivity.

\textbf{Lemma 495.2 (prime-to-$p$ characters in the family).}
Every irreducible complex representation of $A\rtimes P$ has degree a
power of $p$. Its prime-to-$p$ irreducible characters are exactly
\begin{equation}\label{ra495:linearlabels}
 \chi_{\lambda,\mu}(a u)=\lambda(a)\mu(u),
 \qquad \lambda\in\widehat A^{P},\quad\mu\in\widehat P.
\end{equation}

\emph{Proof.}
Let $V$ be an irreducible representation. Since the finite-order commuting
operators from $A$ are simultaneously diagonalizable over $\mathbb C$,
decompose $V$ into its $A$-weight spaces $V_\lambda$. The group $P$
permutes those spaces. The sum over any one orbit is invariant under both
$A$ and $P$, hence under $G$. Irreducibility implies that the occurring
weights form a single $P$-orbit.

Choose a nonzero weight space $V_\lambda$ and its stabilizer
$P_\lambda$. This is an abelian group, so its commuting finite-order
operators on $V_\lambda$ have a common eigenvector $v\ne0$.
The span of all translates $u v$ for $u\in P$ is $P$-invariant and is
$A$-invariant, because each translate belongs to an $A$-weight space.
It is nonzero and therefore equals $V$. Representatives of distinct
cosets of $P_\lambda$ give vectors in distinct weight spaces, while
elements of the same coset give scalar multiples because $v$ is an
eigenvector for $P_\lambda$. The representative vectors are nonzero
and linearly independent. Thus
\[
 \dim V=[P:P_\lambda],
\]
a power of $p$. This degree is prime to $p$ exactly when
$P_\lambda=P$, in which case it is one. A one-dimensional
representation restricts to a $P$-invariant character $\lambda$ of
$A$ and a character $\mu$ of $P$, yielding
\eqref{ra495:linearlabels}. Conversely, invariance of $\lambda$ makes
that formula multiplicative under the semidirect-product law. Distinct
pairs have distinct restrictions to $A$ or to $P$, so give distinct
characters.\hfill$\square$

The normalizer is particularly simple:
\begin{equation}\label{ra495:normalizer}
 N_G(P)=C\times P.
\end{equation}
To prove this, write an element as $a u$ with $a\in A,u\in P$.
The factor $u$ already normalizes $P$. An element $a\in A$ normalizes
$P$ exactly when $a v a^{-1}\in P$ for all $v\in P$. Its commutator
with $v$ lies in $A$ and must also lie in $P$, so it is the identity.
Thus $a\in C$. Conversely every member of $C$ centralizes $P$.
Their intersection is trivial and $P$ is abelian, giving the displayed
direct product, which itself is abelian.

\textbf{Proposition 495.3 (explicit equivariance in this family).}
For $G=A\rtimes P$ as above, restriction to $N_G(P)$ is a bijection
between the prime-to-$p$ irreducible characters. It commutes with every
automorphism of $\mathbb Q(\zeta_{|G|})$, hence with $\mathcal H_p$.

\emph{Proof.}
All the characters in \eqref{ra495:linearlabels} have degree one, and
all irreducible characters of $C\times P$ have degree one. The restriction
map on their labels is
\[
 (\lambda,\mu)\longmapsto(\lambda|_C,\mu),
\]
which is a bijection by \eqref{ra495:dualrestriction}. For any field
automorphism $\sigma$ and any $x\in N_G(P)$,
\[
 (\chi|_{N_G(P)})^\sigma(x)=\sigma(\chi(x))
                 =(\chi^\sigma)|_{N_G(P)}(x).
\]
The maps therefore commute pointwise. This verifies equivariance of the
same explicit bijection, rather than separately selecting unrelated
bijections for different automorphisms.\hfill$\square$

\paragraph{A worked odd-prime example and exact finite checks.}
Take $p=3$, $A=C_7\times C_{13}$, and $P=\langle t\rangle\cong C_3$
acting by
\[
 t\cdot(a,b)=(2a,b).
\]
Since $2^3\equiv1\pmod7$ and $2\not\equiv1\pmod7$, the action
has order three. Here $C=\{0\}\times C_{13}$ and
$N_G(P)=C_{13}\times C_3$ is a proper subgroup of index seven.
Characters of $A$ have labels $(i,j)\in\mathbb Z/7\times\mathbb Z/13$.
Exactly the $13$ labels with $i=0$ are $P$-invariant; the other $78$
labels form $26$ orbits of size three. Hence there are $39$ degree-one
characters, labelled by $(j,k)\in\mathbb Z/13\times\mathbb Z/3$,
and $26$ degree-three characters. The latter can be constructed on the
three weight lines in each orbit, with $P$ cyclically permuting them;
the distinct weights and transitivity prove irreducibility. The degree
sum check is
\[
 39\cdot1^2+26\cdot3^2=273=|G|.
\]
Only the degree-one characters enter the target.

In this example $m=91$, and $3$ has order six modulo $91$: it has
order six modulo $7$ and order three modulo $13$. Thus $\mathcal H_3$
has order $2\cdot6=12$. A permitted exponent $u$ acts on both sides
of the restriction bijection by
\[
 (j,k)\longmapsto(uj\bmod13,uk\bmod3).
\]
Its number of fixed characters is exactly
\begin{equation}\label{ra495:examplefixed}
 \gcd(u-1,13)\,\gcd(u-1,3).
\end{equation}
The accompanying standard-Python script enumerates the character-label
orbits, all twelve permitted exponents modulo $273$, and all $468$
exponent/character pairs. It also verifies the six-point $C_2^2$ example
by computing the fixed-subgroup counts for all five subgroups. These
calculations use modular integers and exact finite sets, not floating
point approximations to roots of unity. They test the proved explicit
family and the finite-set obstruction, not arbitrary character tables.

\paragraph{Stress test: why raw restriction is not a general answer.}
The family proof depended on its prime-to-$p$ characters all being linear.
That fact cannot be assumed in general. For a concrete check, let
$G=S_4$, $p=2$, and
$P=\langle(1234),(13)\rangle$, a Sylow subgroup isomorphic to the
dihedral group of order eight. It is self-normalizing. Indeed its unique
cyclic subgroup of order four is generated by $(1234)$, so every
normalizer of $P$ normalizes this cyclic subgroup; the latter has
normalizer of order eight in $S_4$, since the six four-cycles generate
three conjugate cyclic subgroups. Thus $N_G(P)=P$.

Consider the standard representation
\[
 W=\{(z_1,z_2,z_3,z_4)\in\mathbb C^4:
                          z_1+z_2+z_3+z_4=0\}.
\]
It is irreducible for $S_4$: any nonzero invariant subspace contains a
vector with two unequal coordinates; subtracting its image under that
transposition produces a nonzero multiple of $e_i-e_j$, whose permutation
orbit spans $W$. Its degree three is prime to two. On restriction to $P$,
however, the line spanned by $(1,-1,1,-1)$ is invariant. The usual
Hermitian inner product is permutation invariant, so its orthogonal
complement in $W$ is also invariant and has dimension two. Hence
$W|_P$ is reducible. Raw restriction does not even land in the required
set of irreducible characters in this example.

This does not challenge the proved prime-two theorem; it disproves the
intermediate assertion that the explicit map in Proposition 495.3 always
works. Nor can one simply choose a constituent without further work:
one must prove uniqueness or supply a choice compatible with all of
$\mathcal H_p$, and prove global injectivity and surjectivity for that
choice. The finite-set mark calculation explains why matching sizes of
some fixed subsets is not a replacement for these obligations.

\paragraph{Outcome and the first remaining gap.}
\textbf{Outcome: UNRESOLVED.} The work proves the exact subgroup-mark
criterion for the specified finite Galois action, gives a six-point
counterexample to replacing subgroup marks by individual automorphism
counts, and constructs a fully equivariant bijection for every group
$A\rtimes P$ with $A$ abelian of prime-to-$p$ order and $P$ an abelian
$p$-group. The latter is a rederived solvable-family benchmark, not a
claim to a new frontier result.

For arbitrary finite $G$, the first missing statement in the counting
route is \eqref{ra495:targetmarks} for every subgroup
$U\leq\mathcal H_p$. The ordinary count supplies only $U=1$; the
newer $\mathcal H_0$ theorem supplies marks for subgroups contained in
its action subgroup and does not supply all the remaining marks.
For the constructive route, a replacement for raw restriction must
control non-linear characters and their Galois stabilizers simultaneously.

Specific next steps are to compute exact subgroup marks in odd-prime
examples with noncyclic $\mathcal H_p$, to search for a canonical
constituent operation whose Galois stabilizer is preserved, and to examine
where the coprime abelian decomposition fails in a larger semidirect
product. The current argument supplies neither the missing global
fixed-subgroup equality nor a counterexample to it. It also establishes
no additional blockwise refinement.

\subsection{Sources}
\begin{itemize}
\item[S1]L. Ruhstorfer, A. A. Schaeffer Fry, arXiv:2211.14237 (2022); Adv. Math. 483 (2025), 110369.
\url{https://arxiv.org/abs/2211.14237}.
\textit{Formulation source. Consulted 24 September 2026: The full text states the H\_p-equivariant global bijection and distinguishes its blockwise refinement.}
\item[E1] L. Ruhstorfer and A. A. Schaeffer Fry,
\emph{The McKay--Navarro conjecture for the prime 2},
arXiv:2211.14237v2, 20 May 2025, Introduction and Corollary B.
\url{https://arxiv.org/html/2211.14237v2}.
\textit{Accepted author text of [S1]; consulted 27 September 2026.}
\item[E2] L. Ruhstorfer and A. A. Schaeffer Fry,
\emph{The Isaacs--Navarro Galois Conjecture},
arXiv:2509.02300v1, 2 September 2025, Introduction and Theorem A.
\url{https://arxiv.org/html/2509.02300v1}.
\textit{Recent primary manuscript; only its stated smaller-subgroup theorem
is used as frontier information. Consulted 27 September 2026.}
\item[E3] University of Denver institutional record for
L. Ruhstorfer and A. A. Schaeffer Fry,
\emph{The McKay--Navarro Conjecture for the Prime 2},
Advances in Mathematics 477 (2025), 110369,
DOI 10.1016/j.aim.2025.110369.
\url{https://digitalcommons.du.edu/math_faculty/58/}.
\textit{Primary institutional publication record used for the volume audit;
consulted 27 September 2026.}
\end{itemize}
\end{document}
